\documentclass[10pt,a4paper]{amsart}
\usepackage[margin=19mm]{geometry}
\usepackage{amsmath,amssymb,mathtools}
\usepackage[scr=rsfs,cal=euler]{mathalfa}
\usepackage{xfrac}
\usepackage[unicode,hidelinks,bookmarks=false]{hyperref}
\newcommand{\ie}{i.e.\ }
\newcommand{\cf}{cf.\ }
\newcommand{\resp}{respectively\ }
\newcommand{\Subsection}{\subsection}
\providecommand{\nobreakdash}{\nobreak\hbox{-}}
\setlength{\emergencystretch}{2em}
\multlinegap0pt
\usepackage{amssymb}
\usepackage{enumerate}
\usepackage{tikz}
\usepackage{tcolorbox}
\newtheorem{theorem}{Theorem}
\renewcommand\le{\leqslant}
\title{On applications of the clique-adjacency polynomial\\ to arbitrary finite graphs}
\author{Jake Rigg, John Bamberg}
\date{Source: arXiv:2606.19820v1 (CC BY 4.0)}
\begin{document}
\maketitle

\noindent\emph{Source and licence.} On applications of the clique-adjacency polynomial\\ to arbitrary finite graphs. Version arXiv:2606.19820v1 (CC BY 4.0), distributed by arXiv under the Creative Commons Attribution 4.0 International licence, http://creativecommons.org/licenses/by/4.0/, as recorded in the arXiv licence metadata for this exact version and shown on the abstract page https://arxiv.org/abs/2606.19820v1. Full licence notice: \texttt{licenses/arxiv-2606.19820-CC-BY-4.0.txt}.\par\smallskip

\noindent\textit{Editorial scope.} Counted unit: the article's theorem theorem (source0.tex lines 617--619). The article's complete original proof of it is delivered with it (source0.tex lines 621--636, one \texttt{proof} environment of 16 lines). No mathematics is written, repaired or re-derived: the statement and the proof are the article's own lines, in the article's own notation, and no retained line is substituted or re-flowed.  No further source-local context is needed: every cross-reference of every retained line resolves inside the two delivered spans.  Nothing was omitted from any retained span, so the package carries no omission record.  No retained line contains a citation, so the wrapper carries no bibliography and no bibliography entry.  No retained line contains a figure environment, an image-inclusion command or drawing code, so no image is delivered and the package declares no asset dependency.  Editorial formatting only, in addition: an AMS \texttt{amsart} preamble that restates the article's own declarations and macros used by the retained lines, namely the environments \texttt{theorem} on separate counters, together with the standard packages those lines need.  The retained environments print in this document as Theorem 1 in order of appearance, whereas the article prints the same items with the numbering of its own full document.  The complete native source of the article as distributed in the arXiv e-print for this exact version is bundled unchanged as \texttt{sources/203176/source0.tex}.
	\begin{theorem}
	For a graph $\Gamma$, clique of size $s$ and nested $s$-restricted edge-degree sequence \[\mathcal{D}_s(u_1), \mathcal{D}_s(u_2), \dots, \mathcal{D}_s(u_v)\] we can produce an upper bound $\lambda^*_s$ on $\lambda$ such that \[\lambda_s^* = \frac{\sum_{i=1}^s \sum \mathcal{D}_s(u_i)}{s(s-1)}.\]
	\end{theorem}

	\begin{proof}
	Consider the subsequence of $\mathcal{D}_s(\Gamma)$ of the form $\mathcal{D}_s(u_{1}), \dots, \mathcal{D}_s(u_s)$.  
The set of vertices $\{u_{1}, \dots, u_s\}$ correspond to the subset of $s$ vertices whose $(s-1)$ largest edge-degrees together give the maximum possible total among all $s$-subsets of vertices. Thus, for any $s$-clique $C$, we have that 
    \[  
    \sum_{u\in C}\sum \mathcal{D}_s(u)  
    \le  
    \sum_{i=1}^s\sum \mathcal{D}_s(u_i).  
    \]
    
    This means we have the bound
\[
  \lambda \le 
  \frac{\sum_{i=1}^{s} \sum \mathcal{D}_s(u_i)}{s(s-1)},
\]
and take the right-hand side to be $\lambda^*_s$.
	\end{proof}

\end{document}
